\subsection{Гранаты и лягушки}
%%%%%%%%%%%%%%%%%%%%%%%%%%%%%%%%%%%%%%%%%%%%%%%%%%%%
%%%%%%%%%%%%%%%%%%%%%%%%%%%%%%%%%%%%%%%%%%%%%%%%%%%%
 \z{Артиллерист стреляет из пушки ядром массы $m$ так, чтобы оно упало в неприятельском лагере. На вылетевшее из пушки ядро садится барон Мюнхгаузен, масса которого $5m$. Какую часть пути до неприятельского лагеря ему придется идти пешком?}
%
{$\frac{S}{L} = \frac{35}{36}$}
%%%%%%%%%%%%%%%%%%%%%%%%%%%%%%%%%%%%%%%%%%%%%%%%%%%%
%%%%%%%%%%%%%%%%%%%%%%%%%%%%%%%%%%%%%%%%%%%%%%%%%%%% 
 \z{Неподвижный снаряд разорвался на четыре осколка. Осколки массами $m_1 = 3$ кг, $m_2 = 2$ кг, $m_3 = 4$ кг полетели соответственно со скоростями $v_1 = 200$ м/с вертикально вверх, $v_2 = 150$ м/с горизонтально на север и $v_3 = 100$ м/с горизонтально на восток. Под каким углом к горизонту полетел четвёртый осколок?}
%
{$\varphi \approx 50\degree$}
%%%%%%%%%%%%%%%%%%%%%%%%%%%%%%%%%%%%%%%%%%%%%%%%%%%% 
 \z{Снаряд летит по параболе и разрывается в верхней точке траектории на два одинаковых осколка. Первый осколок падает вертикально вниз, второй — на расстоянии $s$ по горизонтали от места разрыва. Найдите скорость снаряда $v$ перед разрывом, если известно, что разрыв произошел на высоте $h$ и время падения первого осколка равно $\tau$. Трением пренебречь.}
%
{$v = \frac{s g \tau}{4h}$}
%%%%%%%%%%%%%%%%%%%%%%%%%%%%%%%%%%%%%%%%%%%%%%%%%%%%
\z{Лягушка массы $m_1$ сидит на конце доски массы $m_2$ и длины $L$. Доска находится на поверхности пруда. Лягушка совершает прыжок вдоль доски так, что ее начальная скорость $v_0$ (относительно воды) составляет с горизонтом угол $\alpha$ и оказывается на другом конце доски. Определите $v_0$. Сопротивлением воды пренебречь.}
%
{$v_0 = \sqrt{\frac{gL}{(1 + \frac{m_1}{m_2}) \sin 2\alpha}}$}
%%%%%%%%%%%%%%%%%%%%%%%%%%%%%%%%%%%%%%%%%%%%%%%%%%%%	
 \z[{\begin{wrapfigure}[5]{r}{0.3\textwidth}%
		\vspace{-12pt}%
		\begin{flushright}%
 			\includegraphics[width=0.3\textwidth]{6.3.5}%
		\end{flushright}%
	\end{wrapfigure}%
	}]%
	{Сани с седоком и собакой общей массой $M$ съезжают с постоянной скоростью $v_0$ с горы, имеющей уклон $\alpha$ ($\cos\alpha = \frac{6}{7}$). Собака массой $m$ спрыгивает с саней вперёд по ходу их движения и приземляется на склон, имея скорость $v$, направленную под углом $\beta$ к горизонту ($\cos\beta = \frac{3}{7}$). Сани после этого продолжают двигаться по горе вниз. Найти скорость саней с седоком после прыжка собаки.}
 %
{$u = \frac{M v_0 \cos\alpha - m v \cos\beta}{(M - m) \cos\alpha} = \frac{2M v_0 - m v}{2(M - m)}$}
